\chapter[Thème 1 $-$ Convexité]{Thème 1 : Modèles définis par une fonction \\Convexité d'une fonction}
\thispagestyle{fancy}

\section{Dérivée seconde}

\begin{bdefi}
	Soit une fonction \(f\) dérivable sur un intervalle \(I\) dont la dérivée \(f^{\prime}\) est dérivable sur \(I\).\\
	\phantom{$\dfrac{1}{2}$}\pointilles[\textwidth-0.5cm]
	
	\phantom{$\dfrac{1}{2}$}\pointilles[\textwidth-0.5cm]
\end{bdefi}

\begin{bmethodet}{Calculer la dérivée seconde d'une fonction}
	\grasInMet{Enoncé :}
	
	Calculer la dérivée seconde de chacune des fonctions \(f\) et \(g\) définies par :
	\begin{enumMet}
		\item \(f(x)=3 x^{3}-5 x^{2}+1\)
		\item \(g(x)=x e^{x}\)
	\end{enumMet}
	
	\grasInMet{Réponse :}
	
\end{bmethodet}


\section{Fonction convexe et fonction concave}

\subsection{Définitions avec les cordes}

\begin{bdefi}
	\phantom{$\dfrac{1}{2}$}\pointilles[\textwidth-0.5cm]
	\begin{center}
		\includegraphics[scale=0.25]{Img-Cours-01}
	\end{center}
\end{bdefi}


\begin{bdefis}
	Soit une fonction \(f\) définie sur un intervalle \(I\).
	\begin{itemDef}
		\item \phantom{$\dfrac{1}{2}$}\pointilles[\textwidth-0.5cm]
		
		\phantom{$\dfrac{1}{2}$}\pointilles[\textwidth-0.5cm]
		\item \phantom{$\dfrac{1}{2}$}\pointilles[\textwidth-0.5cm]
		
		\phantom{$\dfrac{1}{2}$}\pointilles[\textwidth-0.5cm]
	\end{itemDef}
\end{bdefis}

\begin{minipage}[c]{0.48\linewidth}
	\begin{center}
		\includegraphics[scale=0.25]{Img-Cours-02}
		
		Fonction convexe
	\end{center}
	
\end{minipage} \hfill
\begin{minipage}[c]{0.48\linewidth}
	\begin{center}
		\includegraphics[scale=0.25]{Img-Cours-03}
		
		Fonction concave
	\end{center}
\end{minipage}

\subsection{Définitions avec les tangentes}

\begin{bdefis}
	Soit une fonction \(f\) définie sur un intervalle \(I\).
	\begin{itemDef}
		\item \phantom{$\dfrac{1}{2}$}\pointilles[\textwidth-0.5cm]
		
		\phantom{$\dfrac{1}{2}$}\pointilles[\textwidth-0.5cm]
		\item \phantom{$\dfrac{1}{2}$}\pointilles[\textwidth-0.5cm]
		
		\phantom{$\dfrac{1}{2}$}\pointilles[\textwidth-0.5cm]
	\end{itemDef}
\end{bdefis}

\begin{minipage}[c]{0.48\linewidth}
	\begin{center}
		\includegraphics[scale=0.25]{Img-Cours-04}
		%		\captionsetup{labelformat=empty}
		
		Fonction convexe
	\end{center}
\end{minipage} \hfill
\begin{minipage}[c]{0.48\linewidth}
	\begin{center}
		\includegraphics[scale=0.25]{Img-Cours-05}
		%		\captionsetup{labelformat=empty}
		
		Fonction concave
	\end{center}
\end{minipage}

\begin{bmethodet}{Reconnaître graphiquement la convexité}
	\grasInMet{Enoncé :}
	
	Reconnaître graphiquement la convexité des deux fonctions représentées sur l'intervalle [-3; 5].
	
	\begin{enumMet}
		\begin{minipage}[c]{0.48\linewidth}
			\item \ 
			
			\includegraphics[scale=0.25]{Img-Cours-06}
		\end{minipage} \hfill
		\begin{minipage}[c]{0.48\linewidth}
			\item \ 
			
			\includegraphics[scale=0.25]{Img-Cours-07}
		\end{minipage}
	\end{enumMet}
	
	\grasInMet{Correction : }
	\begin{enumMet}
		\begin{minipage}[c]{0.48\linewidth}
			\item La fonction est concave. Sa courbe est en effet entièrement située en dessous de chacune de ses tangentes.\\
			\includegraphics[scale=0.25]{Img-Cours-08}
		\end{minipage} \hfill
		\begin{minipage}[c]{0.48\linewidth}
			\item La fonction est d'abord convexe puis concave.\\
			\includegraphics[scale=0.25]{Img-Cours-09}
		\end{minipage}
	\end{enumMet}
	\begin{brmq}
		On aurait pu obtenir ses résultats en utilisant les cordes.
	\end{brmq}
\end{bmethodet}

\subsection{Propriétés}

\begin{bprops}
	\begin{itemProp}
		\item \phantom{$\dfrac{1}{2}$}\pointilles[\textwidth-0.5cm]
		\item \phantom{$\dfrac{1}{2}$}\pointilles[\textwidth-0.5cm]
		\item \phantom{$\dfrac{1}{2}$}\pointilles[\textwidth-0.5cm]
		\item \phantom{$\dfrac{1}{2}$}\pointilles[\textwidth-0.5cm]
	\end{itemProp}
\end{bprops}

\begin{bprop}
	Soit une fonction \(f\) définie et deux fois dérivable sur un intervalle \(I\).
	
	\begin{itemProp}
		\item \phantom{$\dfrac{1}{2}$}\pointilles[\textwidth-0.5cm]
		
		\phantom{$\dfrac{1}{2}$}\pointilles[\textwidth-0.5cm]
		\item \phantom{$\dfrac{1}{2}$}\pointilles[\textwidth-0.5cm]
		
		\phantom{$\dfrac{1}{2}$}\pointilles[\textwidth-0.5cm]
	\end{itemProp}
\end{bprop}		

\begin{brmq}
	Dans la pratique, pour étudier la convexité d'une fonction, on détermine le signe de la dérivée seconde.
\end{brmq}

\begin{bmethodet}{Étudier la convexité d'une fonction}
	\grasInMet{Enoncé :}
	
	Soit la fonction \(f\) définie sur \(\mathbb{R}\) par \(f(x)=\dfrac{1}{3} x^{3}-9 x^{2}+4\).\\
	Étudier la convexité de la fonction \(f\).
	
	\grasInMet{Correction : }
	
\end{bmethodet}

\begin{bdefi}
	\begin{minipage}[c]{0.48\linewidth}
		Soit une fonction \(f\) dérivable sur un intervalle \(I\).\\
		Un point d'inflexion est un point où la courbe traverse sa tangente.\\
	\end{minipage} \hfill
	\begin{minipage}[c]{0.48\linewidth}
		\includegraphics[scale=0.25]{Img-Cours-10}
	\end{minipage}
\end{bdefi}

\begin{bprop}
	Au point d'inflexion, la fonction change de convexité.
\end{bprop}

\begin{bexemple}
	
	\begin{minipage}[c]{0.48\linewidth}
		On considère la fonction cube \(x \mapsto x^{3}\).\\
		La tangente en 0 est l'axe des abscisses.\\
		Pour \(x \leq 0\), la courbe est en dessous de sa tangente.\\
		Pour \(x \geq 0\), la courbe est au-dessus de sa tangente.\\
		L'origine est donc le point d'inflexion de la courbe de la fonction cube.\\
		La tangente à la courbe traverse donc la courbe en ce point.\\
	\end{minipage} \hfill
	\begin{minipage}[c]{0.48\linewidth}
		\includegraphics[scale=0.25]{Img-Cours-11}
	\end{minipage}
\end{bexemple}

\begin{bmethodet}{Reconnaître graphiquement un point d'inflexion}
	\grasInMet{Enoncé :}
	
	Déterminer graphiquement le point d'inflexion des fonctions représentées ci-dessous.
	
	\begin{enumMet}
		\begin{minipage}[c]{0.48\linewidth}
			\item \ 
			
			\includegraphics[scale=0.2]{Img-Cours-12}
		\end{minipage} \hfill
		\begin{minipage}[c]{0.48\linewidth}
			\item \ 
			
			\includegraphics[scale=0.2]{Img-Cours-13}
		\end{minipage}
	\end{enumMet}
	
	\grasInMet{Correction : }
	\begin{enumMet}
		\begin{minipage}[c]{0.48\linewidth}
			\item La fonction est d'abord concave puis convexe. Le point de coordonnées(0; 1) semble être un point d'inflexion de la courbe.\\
			\includegraphics[scale=0.2]{Img-Cours-14}\\
		\end{minipage} \hfill
		\begin{minipage}[c]{0.48\linewidth}
			\item La fonction est d'abord convexe puis concave. Le point de coordonnées (2 ; 1) semble être un point d'inflexion de la courbe.\\
			\includegraphics[scale=0.2]{Img-Cours-15}
		\end{minipage}
	\end{enumMet}
\end{bmethodet}

\begin{bmethodet}{Étudier la convexité pour résoudre un problème}
	\grasInMet{Enoncé :}
	
	Une entreprise fabrique des clés USB avec un maximum de 10000 par mois. Le coût de fabrication \(C\) (en milliers d'euros) de \(x\) milliers de clés produites s'exprime par : \(C(x)=0,05 x^{3}-1,05 x^{2}+8 x+4\), définie sur l'intervalle [ \(0 ; 10\) ].
	
	\begin{enumMet}
		\item À l'aide de la calculatrice, conjecturer la convexité de la fonction \(C\).
		
		En déduire si la courbe possède un point d'inflexion.
		\item Démontrer ces résultats.
		\item Interpréter les résultats obtenus au regard du contexte de l'exercice.
	\end{enumMet}
	
	
	
	\grasInMet{Correction : }
	
	\begin{enumMet}
		\item La fonction semble concave sur l'intervalle [ 0 ; 7] et convexe sur l'intervalle [7; 10]. La courbe semble posséder un point d'inflexion pour \(x=7\).\\
		\includegraphics[scale=0.65]{Img-Cours-16-2}
	\end{enumMet}
\end{bmethodet}


\begin{bmethodet}{Étudier la convexité pour résoudre un problème}
	\grasInMet{Enoncé :}
	
	Soit \(f\) la fonction définie sur \(\mathbb{R}\) par \(f(x)=x e^{-\frac{x}{2}}\).\\
	\begin{enumMet}
		\item Calculer la dérivée de la fonction \(f\).
		\item Dresser le tableau de variations de la fonction \(f\).
		\item Tracer la courbe représentative de la fonction \(f\) en s'aidant de la calculatrice.
		\item Déterminer une valeur approchée de l'abscisse du point d'inflexion à la courbe.
		\item Démontrer que \(f^{\prime \prime}(x)=\left(\frac{x}{4}-1\right) e^{-\frac{x}{2}}\).
		\item En déduire l'abscisse du point d'inflexion.
	\end{enumMet}
	
	
	
	\grasInMet{Correction : }
\end{bmethodet}
